% language: swedish
% figure: fig1 0.32 0.65 0.58 0.69
\section{Repetition}
\subsection{Komplex deriverbarhet}
$f : G \to \mathbb{C}$ \emph{komplext deriverbar} i $z_0 \in G$ om $f'(z_0) = \displaystyle\lim_{h \to 0} \frac{f(z_0 + h) - f(z_0)}{h}$ existerar. Om $G$ öppen och $f$ komplext deriverbar i alla punkter i $G$ då säger vi att $f$ är holomorf i $G$.

\begin{theorem}[Cauchy-Riemanns ekv.\ (Sats 2.13)]
a) om $f$ holomorf, $u = \operatorname{Re} f$, $v = \operatorname{Im} f$, då löser $u, v$ CR ekv.
\[ \begin{cases} \dfrac{du}{dx} = \dfrac{dv}{dy} \\[2ex] \dfrac{du}{dy} = -\dfrac{dv}{dx} \end{cases} \qquad \text{alt.} \quad \frac{df}{dx} = -i\frac{df}{dy} \ (= f') \]
b) om $f \in C^1$ och dess partiella derivator löser CR:s ekv då är $f$ holomorf.
\end{theorem}

\begin{example}
$e^z = e^x(\cos y + i\sin y)$, såg $e^z$ holomorf i $\mathbb{C}$ (hel) mha CR:s ekv.
\end{example}

\begin{theorem}[Sats 2.17]
Om $f$ holomorf i område $G$ (dvs.\ öppen \& sammanhängande) och $f' \equiv 0$ i $G$, då är $f$ konstant
\end{theorem}

\begin{example}
Antag att $f$ holomorf och reellvärd i $G$. Visa att $f$ konstant i $G$. \hfill \textcolor{magenta!80!black}{område}

\textbf{Lösning:} Låt $u = \operatorname{Re} f$, $v = \operatorname{Im} f$, enligt antagande $v \equiv 0$.
\[ \frac{du}{dx} \overset{\textcolor{magenta!80!black}{\text{CR}}}{=} \frac{dv}{dy} \equiv 0 \qquad f' \overset{\textcolor{magenta!80!black}{\text{CR}}}{=} \frac{df}{dx} = \frac{du}{dx} + i\frac{dv}{dx} = 0 \underset{\textcolor{magenta!80!black}{\text{Sats}}}{\Longrightarrow} f \text{ konstant} \]
\end{example}

\subsection{Konforma avbildningar}
\begin{definition}
En komplex fkn $f : G \to \mathbb{C}$ är \emph{konform} om den bevarar vinkeln mellan kurvor i $G$
\notefigure{fig1}{}
\end{definition}

\textbf{Proposition 2.11:} Om $f$ holomorf i $G$, $f' \neq 0$ i $G$, då är $f$ \emph{konform}

\textbf{Lemma:} $f$ holomorf, $\gamma$ kurva $(f \circ \gamma)'(t) = f'(\gamma(t))\gamma'(t)$

\subsection{Möbiusavbildningar}
$M(z) = \dfrac{az + b}{cz + d}$ kallas Mavb om $ad - bc \neq 0$. $M'(z) = \dfrac{ad - bc}{(cz + d)^2} \neq 0$

dvs.\ Mavb är konform. $\hat{\mathbb{C}} = \mathbb{C} \cup \{\infty\}$ \quad $\dfrac{az + b}{cz + d}$, $\hat{\mathbb{C}} \to \mathbb{C}$, $-\dfrac{d}{c} \to \infty$, $\infty \to \dfrac{a}{c}$

Får bijektion.

\textbf{Sats 3.4:} Mavb avbildar cirklar \& linjer på cirklar eller linjer.
